\chapter{La Campaña en 4 Fases: La Lotería Azul}

La aventura está estructurada en cuatro fases correlativas, pensadas para resolverse en 3 a 4 sesiones de juego de 3 a 4 horas cada una.

\section{Fase 1: El despertar en el asfalto (La Investigación)}

\subsection{Escena 1.1: El secreto en la trastienda}
La campaña arranca en el reservado de la Tasca «Zulo». Mientras el sirimiri repiquetea contra el cristal, la familia de Ane Garate suplica ayuda a los personajes.
\begin{itemize}
  \item \lead{El examen de Ane:} Si el Celador o la Biohacker examinan la nuca de Ane, comprueban que el puerto no es un conector de memoria ordinario: es un bus de acoplamiento neuronal de alta velocidad con cuatro líneas de fibra de cuarzo atornilladas a la base del cráneo.
  \item \lead{La baliza activa:} El puerto emite un parpadeo de LED azul intermitente. La Biohacker descubre que está radiando una señal de telemetría médica encriptada. Cada diez minutos de indecisión marca 1 segmento en el \textbf{Reloj de Limpieza de los Auditores}.
\end{itemize}

\subsection{Escena 1.2: Allanamiento en el Dispensario Periférico}
Para comprender qué le han inyectado a Ane y estabilizar su fiebre, el grupo necesita su historial clínico original o una dosis de Neuro-Ciclotox.
\begin{itemize}
  \item \lead{Infiltración en el archivo:} El Burócrata Arrepentido puede emplear su movimiento «Sello de goma» para engañar al administrativo de guardia, o el Sabueso puede forzar una salida de incendios en el callejón trasero.
  \item \lead{El hallazgo revelador:} Los archivos de los ganadores del Programa Azul no contienen certificados de traslado a la Cúpula Dorada, sino actas de defunción falsificadas por parada cardiorrespiratoria selladas por la \textbf{División de Proyectos Especiales IZAR}.
\end{itemize}

\subsection{Escena 1.3: Emboscada en la estación del tren mag-lev}
Al intentar trasladar a Ane hacia un refugio seguro, dos Auditores de gabardina interceptan al grupo en las escaleras mecánicas de la Línea 4.
\begin{itemize}
  \item \lead{Conflicto dinámico:} Los Auditores no disparan; activan un bloqueador sónico que aturde a la multitud y avanzan con inyectores de sedante. El grupo debe decidir entre el choque frontal, provocar una estampida de viajeros o saltar a las vías magnéticas antes de que llegue el convoy.
  \item \lead{Objetivo cumplido:} Confirmar la conspiración corporativa y obtener el nombre clave: \textbf{Proyecto IZAR}.
\end{itemize}

\section{Fase 2: El subsuelo y los Descartados (El Descenso)}

\neofigure{09_descenso_subsuelo_alcantarillas.jpg}{Descenso por los colectores de la Ría hacia el refugio subterráneo de los Descartados.}

\subsection{Escena 2.1: La marcha por las cloacas ácidas}
La superficie del Distrito 7 está infestada de patrullas. La única vía para salvar a Ane es descender por las bocas del colector pluvial hacia los cimientos inundados de la vieja Ría.
\begin{itemize}
  \item \lead{Desafío ambiental:} El lodo químico arrastra vapores corrosivos. Los jugadores deben coordinar sus paraguas y máscaras de filtración. Un fallo en tirada de \textbf{Asustado} deteriora el equipo y añade 1 punto de Zama por claustrofobia.
\end{itemize}

\subsection{Escena 2.2: La comuna de los cerebros rotos}
En las bóvedas de hormigón de un antiguo muelle de carga subterráneo, el grupo descubre el campamento de los Descartados. Decenas de antiguos ciudadanos brillantes sobreviven entre hogueras de residuos plásticos, tarareando al unísono melodías sintéticas en trance.
\begin{itemize}
  \item \lead{Encuentro con el Dr. Bittor Albéniz:} El científico desertor acepta conectar el puerto de Ane a su decodificador artesanal. Tras varios minutos de tensión eléctrica donde las pantallas CRT emiten chispazos, Albéniz extrae el mapa de acceso al búnker del Proyecto IZAR y la tarjeta criptográfica del jefe de seguridad.
  \item \lead{El asalto de Limpieza Sanitaria:} Dos furgonetas blancas irrumpen por una esclusa de servicio. Los paramédicos comienzan a gasear el campamento para eliminar testigos. Los PJs deben elegir entre huir de inmediato hacia el montacargas o arriesgar sus vidas para salvar a los Descartados indefensos.
\end{itemize}

\section{Fase 3: El ascenso a la Cúpula (La Infiltración)}

\neofigure{10_infiltracion_cupula_blanca.jpg}{El violento contraste: pasillos asépticos de mármol blanco en la Torre Orekak.}

\subsection{Escena 3.1: El choque estético del Nivel Superior}
A través del montacargas industrial, los personajes emergen en el vestíbulo de la Torre Orekak. El contraste sensorial es brutal: del hedor a azufre y la oscuridad del subsuelo pasan a suelos de mármol sintético inmaculado, luz blanca relajante y música ambiental esterilizada.
\begin{itemize}
  \item \lead{Mecánica de Alarma de la Cúpula (0/6):} Los personajes están cubiertos de barro y hollín. Deben conseguir uniformes de mantenimiento, batas médicas o trajes de oficina en menos de diez minutos antes de que los sensores biométricos de techo detecten «anomalías estéticas».
\end{itemize}

\subsection{Escena 3.2: El descenso al Búnker Subterráneo -4}
El laboratorio no está en la cúspide de la torre, sino en sus cimientos más blindados.
\begin{itemize}
  \item \lead{Superar los filtros de seguridad:} El Burócrata Arrepentido y el Sabueso deben coordinarse: el primero distrae al oficial de guardia con auditorías inventadas mientras la Biohacker puentea el escáner de retina del ascensor de alta velocidad.
\end{itemize}

\section{Fase 4: El Evento de Resonancia Azul (El Clímax)}

\neofigure{11_evento_resonancia_azul.jpg}{La sala del mainframe colapsando: fulgor azul eléctrico y tormenta psíquica desatada.}

\subsection{Escena 4.1: El santuario de la pesadilla biológica}
Las puertas acorazadas se abren a una sala abovedada de treinta metros de altura. En el centro, sumergidos en cilindros de cristal llenos de líquido amniótico azul fosforescente, flotan dieciséis adultos desnudos. Decenas de cables de fibra óptica nacen de sus nucas y convergen en un cilindro negro central: el procesador \textbf{ADUR-OS v1.0}.

En la pasarela de control, el \textbf{Dr. Victor Kessler} supervisa el encendido final junto a cuatro guardias armados.

\subsection{Escena 4.2: La catástrofe en directo}
En el instante en que los personajes irrumpen o intentan sabotear la consola, el sistema alcanza el 80\% de sincronización y se quiebra.
\begin{itemize}
  \item \lead{La tormenta de Resonancia:} Las mentes de los dieciséis sujetos entran en pánico colectivo. Se desata una onda de choque telequinética y electromagnética. Los cristales de los monitores CRT estallan en una lluvia de esquirlas, las luces de emergencia giran en rojo sangre y los personajes sienten en sus propios cerebros los gritos desgarradores de los sujetos conectados.
  \item \lead{Ataque de Zama masivo:} Cada turno que los personajes permanezcan en la sala, deben superar una tirada de emoción o sufrir 1 punto de Zama directo por la intrusión psíquica.
\end{itemize}

\subsection{Escena 4.3: Los tres nodos de desactivación}
Para detener la catástrofe antes de que sus cerebros hiervan o lleguen las tropas de élite, el grupo debe neutralizar tres subsistemas críticos:
\begin{enumerate}
  \item \lead{El Bus Sináptico Central (La Biohacker):} Puentear el rack de servidores tirando + \textbf{Culpable} para interrumpir el flujo de datos sin freír las cortezas de los pacientes.
  \item \lead{Las Válvulas de Nitrógeno Líquido (El Sabueso o Familiar):} Tirar + \textbf{Enfadado} para forzar manualmente los volantes oxidados de refrigeración y ahogar el sobrecalentamiento del reactor biológico.
  \item \lead{El Desacoplo de los Sujetos (El Celador o Burócrata):} Tirar + \textbf{Afligido} o + \textbf{Pacífico} para drenar los tanques y aplicar antídotos de emergencia a los supervivientes.
\end{enumerate}

\subsection{El Desenlace y la Encrucijada Moral}
Con el colapso del sistema y Kessler derrotado o acorralado, los jugadores se enfrentan a la decisión final:
\begin{itemize}
  \item \textbf{Opción A: Divulgar la verdad.} Conectar la cinta con los datos del Proyecto IZAR al transmisor de propaganda de la torre. La verdad sobre la Lotería Azul saldrá en todas las pantallas de La Margen, desatando huelgas salvajes y disturbios masivos.
  \item \textbf{Opción B: El rescate silencioso.} Evacuar a los pocos supervivientes por los túneles de drenaje y volar el complejo con una sobrecarga eléctrica, borrando las pruebas para proteger a sus familias de represalias corporativas inmediatas.
\end{itemize}
